\section{Worked Example 1: 一个 toy perplexity 计算}

为了把 perplexity 讲透，我们做一个非常小的例子。假设一个模型对两个 token 的条件概率分别是：
\[
p(x_1)=0.5,\qquad p(x_2|x_1)=0.25.
\]
那么这两个 token 的联合概率是
\[
p(D)=0.5\times 0.25=0.125.
\]
如果数据集长度是 2，那么 perplexity 就是
\[
\mathrm{PPL}(D)=\left(\frac{1}{0.125}\right)^{1/2}=\sqrt{8}\approx 2.83.
\]
这意味着，模型平均上相当于在“每一步从大约 2.83 个等可能选择中猜一个”。

\begin{knowledgebox}{这个例子的意义}
perplexity 不是抽象术语，它就是把序列概率转成更直观的“平均分叉度”。数值越小，模型对真实文本越不惊讶。
\end{knowledgebox}

\begin{warningbox}{toy 例子和真实评估的差异}
真实场景里，模型面对的不是两个 token，而是数百到数万 token；概率也不是均匀的，所以 perplexity 的可解释性会和任务难度、上下文长度、tokenization 共同耦合。
\end{warningbox}

